% Part: incompleteness
% Chapter: theories-computability
% Section: interpretability

\documentclass[../../../include/open-logic-section]{subfiles}

\begin{document}

\olfileid{inc}{tcp}{itp}

\olsection{النظريات التي تقبل تأويل $\Th{Q}$ غير قابلة للقرار}

ويمكن تجاوز هذه النتائج إلى ما هو أقوى. فتأويل اللغة
$\Lang{L_1}$ في لغة أخرى $\Lang{L_2}$ معناه، على وجه غير صوري،
أن نعرّف كون $\Lang{L_1}$ ورموز علاقاتها ودوالها باستعمال
!!{formula}s من $\Lang{L_2}$. ويمكن صوغ هذا المعنى في تعريف دقيق،
لكننا لا نفصّل ذلك هنا.

\begin{thm}
  لتكن $\Th{T}$ نظرية في لغة تقبل تأويل لغة الحساب بحيث تبقى
  $\Th{T}$ متسقة مع تأويل $\Th{Q}$. فلا تكون $\Th{T}$ قابلة
  للقرار. فإن كانت $\Th{T}$ تثبت تأويل بديهيات $\Th{Q}$، امتنع
  كذلك أن يكون لـ$\Th{T}$ امتداد متسق قابل للقرار.
\end{thm}

نحصل على البرهان بتغيير يسير في برهان المبرهنة السابقة؛ فوجود
مثال مضاد يفضي إلى فصل $\Th{Q}$ عن $\Th{\bar Q}$.
وللتطبيق نأخذ $\Th{ZFC}$، نظرية مجموعات زرميلو--فرينكل مع بديهية
الاختيار، أساسًا بديهيًا يكفي لبناء جانب واسع من الرياضيات المعتادة.
ففي $\Th{ZFC}$ يمكن تعريف الأعداد الطبيعية، وبذلك التأويل تصدق
بديهيات $\Th{Q}$. ومن هنا نستنتج:

\begin{cor}
يمتنع وجود امتداد متسق قابل للقرار للنظرية $\Th{ZFC}$.
\end{cor}

\begin{cor}
لا تجتمع خصائص التمام والاتساق وكون النظرية !!{axiomatizable}
حسابيًا في امتداد لـ$\Th{ZFC}$.
\end{cor}

ولغة $\Th{ZFC}$ لا تشتمل إلا على علاقة ثنائية واحدة، هي $\in$،
بل لا يلزم استعمال الهوية نفسها. فينتج:

\begin{cor}
كل لغة تشتمل على رمز علاقة ثنائية يكون منطق الرتبة الأولى فيها
غير قابل للقرار.
\end{cor}

ويكفي أيضًا أن تشتمل اللغة على رمزين لدالتين أحاديتين، إذ بهما
يمكن محاكاة رمز العلاقة الثنائية. وهذه الحدود دقيقة: فمنطق الرتبة
الأولى يقبل القرار إذا اقتصرت اللغة على رموز علاقات \emph{أحادية}
وعلى رمز دالة \emph{أحادي} واحد على الأكثر.

ويحسن أن نذكر هنا نتيجة أخرى. قد علمنا أن جمل اللغة ذات الرموز
$\Obj 0$، و$'$، و$+$، و$\times$، و$<$ الصادقة في النموذج القياسي
تؤلف مجموعة غير قابلة للقرار. ويمكن تعريف $<$ ببقية الرموز،
ثم تعريف $+$ باستعمال $\times$ و$'$؛ فتبقى مجموعة الجمل الصادقة
غير قابلة للقرار حتى في اللغة ذات الرموز $\Obj 0$، و$'$، و$\times$.
وأما اللغة ذات الرموز $\Obj 0$، و$'$، و$+$، فقد أثبت بريسبورغر أن
مجموعة جملها الصادقة في النموذج القياسي قابلة للقرار، وإن كان إجراء
القرار لا يصلح عمليًا لشدة كلفته الحسابية.

\end{document}
